\section{Practice Problems}
\label{sec:practice}

The exercises of the slides and further exam-style problems, with complete
solutions; the ones the slides left to the students are solved too. Every number
was recomputed by a small program. The final exam review
(\texttt{reviews/final\_exam\_review/}) holds the lecturer's review questions and
more exam-style questions on every topic.

\providecommand{\enc}[2]{\par{\centering\texttt{#1} $=$ \texttt{#2}\par}}

\subsection{Performance, Cost and Amdahl's Law}

\begin{exercise}[Unix \texttt{time}]
The command prints \texttt{90.7u 12.9s 2:39 65\%}. Interpret it.
\end{exercise}
\begin{solution}
User CPU $=90.7$~s, system CPU $=12.9$~s, elapsed $=2{:}39=159$~s. CPU
time $=103.6$~s, and $103.6/159 = 65\%$ of the elapsed time. The other
35\% is I/O or time spent running other programs.
\end{solution}

\begin{exercise}[Two notions of performance]
Boeing 747: 6.5~h DC--Paris, 610~mph, 470 passengers. Concorde: 3~h,
1350~mph, 132 passengers. Compare latency and throughput.
\end{exercise}
\begin{solution}
Latency: Concorde is $6.5/3 = 1350/610 = 2.2\times$ faster.
Throughput (passenger-miles per hour): Boeing $470\times610=286{,}700$,
Concorde $132\times1350=178{,}200$. Boeing is $286{,}700/178{,}200 =
1.6\times$ higher throughput. Which is ``faster'' depends on the metric.
\end{solution}

\begin{exercise}[Price-performance: train vs.\ plane]
Move 2400 people. Train: 4~h trip, 2400 seats, 300~Rs/person. Plane: 45~min
trip, 300 seats, 3000~Rs/person. Compare time per person and cost-performance.
\end{exercise}
\begin{solution}
Plane needs $2400/300=8$ trips $=6$~h; train needs 4~h. Time per person:
train $4\text{ h}/2400 = 6$~s, plane $6\text{ h}/2400=9$~s, so the train
has 50\% better throughput even though each plane trip is $5\times$
faster. Cost $\times$ time: train $300\times6=1800$, plane
$3000\times9=27{,}000$ Rs$\cdot$s/person: the train is $15\times$ better.
\end{solution}

\begin{exercise}[Speedup]
A program takes 1~s on machine A and 10~s on machine B. Speedup of A over B?
\end{exercise}
\begin{solution}
$\text{Speedup} = \text{ExTime}_B/\text{ExTime}_A = 10/1 = 10$: A is
10 times faster.
\end{solution}

\begin{exercise}[CPU time]
$I = 10{,}000{,}000$ instructions, $\CPI=2.5$, clock 200~MHz. CPU time?
\end{exercise}
\begin{solution}
$C = 1/200\text{ MHz} = 5\times10^{-9}$~s, so
$T = 10^7\times2.5\times5\times10^{-9} = 0.125$~s.
\end{solution}

\begin{exercise}[New compiler and faster CPU]
Same program. A new compiler gives $I'=9{,}500{,}000$ and $\CPI'=3.0$; the
new CPU runs at 300~MHz. Speedup?
\end{exercise}
\begin{solution}
$T' = 9.5\times10^6\times3.0/(300\times10^6) = 0.095$~s.
Speedup $=0.125/0.095 = 1.32$, i.e.\ 32\% faster.
\end{solution}

\begin{exercise}[Instruction classes]
$\CPI_A=1$, $\CPI_B=2$, $\CPI_C=3$. Sequence 1 has counts $(2,1,2)$,
sequence 2 has $(4,1,1)$. Find each CPI. Which is faster?
\end{exercise}
\begin{solution}
Seq.~1: cycles $=2+2+6=10$, $I=5$, $\CPI=2.0$.
Seq.~2: cycles $=4+2+3=9$, $I=6$, $\CPI=1.5$.
Sequence 2 is faster (9 vs.\ 10 cycles) even though it executes \emph{more}
instructions.
\end{solution}

\begin{exercise}[RISC instruction mix]
ALU 50\% (CPI 1), Load 20\% (CPI 5), Store 10\% (CPI 3), Branch 20\% (CPI 2).
Find the CPI and the fraction of time spent in each class.
\end{exercise}
\begin{solution}
$\CPI = 0.5+1.0+0.3+0.4 = 2.2$. Time fractions: ALU $0.5/2.2=23\%$, Load
$1.0/2.2=45\%$, Store $0.3/2.2=14\%$, Branch $0.4/2.2=18\%$.
\end{solution}

\begin{exercise}[$n\%$ faster]
Y takes 15~s, X takes 10~s. How much faster is X?
\end{exercise}
\begin{solution}
$n = 100\,(15-10)/10 = 50\%$.
\end{solution}

\begin{exercise}[Amdahl: floating point]
FP instructions become $2\times$ faster, but only 10\% of the time is FP.
Overall speedup?
\end{exercise}
\begin{solution}
$\text{ExTime}_{new} = (0.9 + 0.1/2)\,\text{ExTime}_{old} =
0.95\,\text{ExTime}_{old}$, so Speedup $=1/0.95 = 1.053$.
\end{solution}

\begin{exercise}[Amdahl: load CPI]
In the RISC mix above, load CPI improves from 5 to 2. Speedup?
\end{exercise}
\begin{solution}
Loads take $F=45\%$ of the time (not 20\%), $S=5/2=2.5$:
$\text{Speedup}=1/(0.55+0.45/2.5) = 1/0.73 = 1.37$.
Check with CPI: new $\CPI = 0.5+0.2\times2+0.3+0.4 = 1.6$, and
$2.2/1.6 = 1.37$.
\end{solution}

\begin{exercise}[How fast must multiply be?]
A program runs in 100~s, of which multiply takes 80~s. How much faster must
multiply be to make the program (a) $4\times$ faster, (b) $5\times$ faster?
\end{exercise}
\begin{solution}
(a) Target 25~s: $25 = 20 + 80/n \Rightarrow n = 16$.
(b) Target 20~s: $20 = 20 + 80/n \Rightarrow 80/n = 0$: impossible. The
unaffected 20~s already equals the target ($1/(1-F) = 1/0.2 = 5$ is the
limit and cannot be reached).
\end{solution}

\begin{exercise}[Multiple enhancements]
$S_1=10, F_1=20\%$; $S_2=15, F_2=15\%$; $S_3=30, F_3=10\%$ (disjoint
parts). Overall speedup?
\end{exercise}
\begin{solution}
$\text{Speedup} = 1/\big(0.55 + 0.02 + 0.01 + 0.0033\big) = 1/0.5833 =
1.71$.
\end{solution}

\begin{exercise}[MIPS can mislead]
Clock 100~MHz; CPI of classes A, B, C $=1,2,3$. Compiler 1 produces
$(5,1,1)$ million instructions; compiler 2 produces $(10,1,1)$ million.
Compare MIPS and CPU time.
\end{exercise}
\begin{solution}
Compiler 1: $\CPI = 10/7 = 1.43$, MIPS $=100/1.43 = 70$, time
$=7\times10^6\times1.43/10^8 = 0.10$~s.
Compiler 2: $\CPI = 15/12 = 1.25$, MIPS $= 80$, time
$=12\times10^6\times1.25/10^8 = 0.15$~s.
Compiler 2 has the higher MIPS but is \emph{slower}.
\end{solution}

\begin{exercise}[Dies per wafer]
Die $0.7$~cm on a side, wafer diameter 30~cm. How many dies?
\end{exercise}
\begin{solution}
$A=0.49\text{ cm}^2$:
$\pi(15)^2/0.49 - \pi\cdot30/\sqrt{0.98} = 1442.6 - 95.2 \approx 1347$ dies.
\end{solution}

\begin{exercise}[Die yield]
Same die, defect density $0.6$/cm$^2$, $\alpha=4$, wafer yield 100\%.
\end{exercise}
\begin{solution}
$(1 + 0.6\times0.49/4)^{-4} = 1.0735^{-4} = 0.75$.
\end{solution}

\begin{exercise}[Weighted means]
P1/P2 take 1/1000~s on A, 10/100~s on B, 20/20~s on C. Rank the machines
under weights $(0.5,0.5)$, $(0.909,0.091)$, $(0.999,0.001)$.
\end{exercise}
\begin{solution}
\begin{tabular}[t]{@{}lccc@{}}
Weights & A & B & C \\
\hline
$.5/.5$     & 500.5 & 55.0  & \textbf{20} \\
$.909/.091$ & 91.91 & \textbf{18.19} & 20 \\
$.999/.001$ & \textbf{2.0} & 10.09 & 20 \\
\end{tabular}

Each weighting picks a different winner (C, then B, then A). 
\end{solution}
\begin{remark}
\slidefix{} The slide prints 91.82 and 18.8 for the middle row; the exact values
are 91.91 and 18.19. The ranking is the same.
\end{remark}

\subsection{Instruction Sets and MIPS}

\begin{exercise}[Encode an R-format instruction]
Encode \texttt{add \$t0, \$s1, \$s2}.
\end{exercise}
\begin{solution}
opcode $0$, rs $=$ \texttt{\$s1} $=17$, rt $=$ \texttt{\$s2} $=18$, rd $=$ \texttt{\$t0} $=8$, shamt $0$, funct
\texttt{0x20}:
\enc{000000 10001 10010 01000 00000 100000}{0x02324020}
\end{solution}

\begin{exercise}[Shifts]
Encode \texttt{sll \$t2,\$s0,8} and \texttt{sra \$t2,\$s0,8}. Why is there no
\texttt{sla}, and how is \texttt{NOT} done?
\end{exercise}
\begin{solution}
Shifts use rt as source and shamt as the amount; rs $=0$.
\texttt{sll}: \enc{000000 00000 10000 01010 01000 000000}{0x00105200}
\texttt{sra}: same with funct \texttt{000011} $=$ \texttt{0x00105203}.
A left arithmetic shift would be identical to \texttt{sll} (zeros enter from
the right either way), so it is not needed. \texttt{NOT} is
\texttt{nor \$t0,\$t1,\$zero}.
\end{solution}

\begin{exercise}[I-format arithmetic and logical]
Encode \texttt{addi \$t0,\$t1,0x0005} and
\texttt{ori \$t0,\$t1,0xFF00}. Why is there no \texttt{subi}?
\end{exercise}
\begin{solution}
rs $=$ \texttt{\$t1} $=9$, rt $=$ \texttt{\$t0} $=8$.
\texttt{addi}: \enc{001000 01001 01000 0000000000000101}{0x21280005}
\texttt{ori}: \enc{001101 01001 01000 1111111100000000}{0x3528FF00}
\texttt{subi} is unnecessary: \texttt{addi} with a negative (sign-extended)
immediate does the same job.
\end{solution}

\begin{exercise}[Loads and stores]
Encode \texttt{lw \$t0,16(\$t1)} and \texttt{sw \$t0,16(\$t1)}.
\end{exercise}
\begin{solution}
rs $=$ base \texttt{\$t1} $=9$, rt $=$ \texttt{\$t0} $=8$, imm $=16$.
\texttt{lw}: \enc{100011 01001 01000 0000000000010000}{0x8D280010}
\texttt{sw}: opcode \texttt{101011} $\Rightarrow$ \texttt{0xAD280010}.
\end{solution}

\begin{exercise}[Branch]
Encode \texttt{beq \$s0,\$s1,256} (offset field $=256$).
\end{exercise}
\begin{solution}
Encoding:
\enc{000100 10000 10001 0000000100000000}{0x12110100}
The target is PC$+4+256\times4$ = PC$+4+1024$.
\end{solution}

\begin{exercise}[Jump]
Encode \texttt{j 0x01000000} as shown on the slide.
\end{exercise}
\begin{solution}
The slide places \texttt{0x01000000} directly in the 26-bit field:
\enc{000010 01000000000000000000000000}{0x09000000}
The real target is $\{\text{PC}{+}4[31{:}28],\ \text{field},\ 00\}$, i.e.\
the field is a \emph{word} address. If you are given a byte target address
instead, divide it by 4 first (divide a byte address by 4).
\end{solution}

\begin{exercise}[Move without MOV]
Move \texttt{\$t0} into \texttt{\$t1} with one instruction.
\end{exercise}
\begin{solution}
\texttt{add \$t1,\$t0,\$zero} (\texttt{or \$t1,\$t0,\$zero} or
\texttt{addi \$t1,\$t0,0} also work).
\end{solution}

\begin{exercise}[Large constant]
Load \texttt{0x12345678} into \texttt{\$t0}.
\end{exercise}
\begin{solution}
\texttt{lui \$t0,0x1234} then \texttt{ori \$t0,\$t0,0x5678}. \texttt{ori}
zero-extends its immediate, so the upper half is not disturbed.
\end{solution}

\begin{exercise}[Atomic swap with LL/SC]
Explain the loop
\texttt{try: add \$t0,\$zero,\$s4; ll \$t1,0(\$s1); sc \$t0,0(\$s1);
beq \$t0,\$zero,try; add \$s4,\$zero,\$t1}.
\end{exercise}
\begin{solution}
\texttt{ll} reads the old memory value into \texttt{\$t1}. \texttt{sc}
tries to store \texttt{\$s4}'s copy; it writes 1 into \texttt{\$t0} on
success, or 0 if the location changed since the \texttt{ll} (another
processor intervened), in which case the loop retries. On success
\texttt{\$s4} receives the old value: memory and \texttt{\$s4} have been
swapped atomically.
\end{solution}

\begin{exercise}[Procedure-call convention]
(1) Who saves the return address, caller or callee? (2) Who saves the
necessary registers? (3) Must a leaf procedure save the return address?
\end{exercise}
\begin{solution}
(1) The caller's \texttt{jal} writes the return address into \texttt{\$ra};
it is the \textbf{callee} that must preserve it (push \texttt{\$ra} on the
stack) if the callee itself calls other procedures, because \texttt{\$ra} is
``preserved across a call''. (2) Both share the work: the \textbf{callee}
saves the \texttt{\$s} registers it uses (callee-saved); the \textbf{caller}
saves any \texttt{\$t}/\texttt{\$a}/\texttt{\$v} registers it still needs
after the call (caller-saved). (3) \textbf{No}: a leaf procedure makes no
\texttt{jal}, so \texttt{\$ra} is never overwritten and it can return with
\texttt{jr \$ra} directly.
\end{solution}

\begin{exercise}[Review quiz]
(1) In MIPS, \texttt{NOT} is done with: a.\ \texttt{NOT}, b.\ \texttt{NOR},
c.\ \texttt{NAND}, d.\ none.
(2) An addressing mode is: a.\ how an operand (or a jump target) is
referenced, b.\ how memory is organized by byte/halfword/word, c.\
endianness, d.\ all of these.
(3) MIPS \texttt{LO}/\texttt{HI} are used for: a.\ low/high addresses, b.\
multiply/divide, c.\ both, d.\ neither.
(4) MIPS stores data as: a.\ little-endian, b.\ big-endian, c.\ both, d.\
no endianness.
\end{exercise}
\begin{solution}
(1) \textbf{b} -- \texttt{nor \$t0,\$t1,\$zero}.
(2) \textbf{a}.
(3) \textbf{b} -- \texttt{mult}/\texttt{div} put results in Hi/Lo
(\texttt{mfhi}/\texttt{mflo} read them).
(4) \textbf{b} -- big-endian, as stated in the lecture (real MIPS chips can
be configured either way, but the expected quiz answer is b).
\end{solution}

\begin{exercise}[C to MIPS]
Which sequences implement \texttt{if (a==0) a=1; else a=0;}?
\begin{enumerate}[label=(\Alph*),leftmargin=*]
  \item \texttt{lw \$t0,a; bne \$t0,\$0,L1; addi \$t0,\$0,1; b L2;}\\
        \texttt{L1: add \$t0,\$0,\$0; L2: sw \$t0,a}
  \item \texttt{lw \$t0,a; beq \$t0,\$0,L1; sw \$0,a; b L2;}\\
        \texttt{L1: addi \$t0,\$0,1; sw \$t0,a; L2:}
  \item \texttt{lw \$t0,a; sltiu \$t0,\$t0,1; sw \$t0,a}
  \item \texttt{lw \$t0,a; addi \$t1,\$0,1; beq \$t0,\$0,L1;}\\
        \texttt{addi \$t1,\$0,0; L1: sw \$t1,a}
  \item All of the above.
\end{enumerate}
\end{exercise}
\begin{solution}
\textbf{E}. (A) jumps to L1 when $a\neq0$ and stores 0, otherwise stores 1.
(B) branches to L1 when $a=0$ and stores 1, otherwise stores \texttt{\$0}.
(C) \texttt{sltiu} sets \texttt{\$t0} $=1$ iff $a<1$ unsigned, i.e.\ iff
$a=0$. (D) presets 1 and overwrites it with 0 only when $a\neq0$.
Caveat: this assumes no branch delay slots (the assembler/simulator
default). With delay slots exposed, the instruction after every branch
always executes: in (A) the \texttt{add} after \texttt{b L2} would zero
\texttt{\$t0} before the store, and in (D) the \texttt{addi \$t1,\$0,0}
would always run, so both would store 0. (B) and (C) would still be
correct.
\end{solution}

\subsection{Pipelining and Hazards}

\begin{exercise}[Laundry]
Wash 30~min, dry 40~min, fold 20~min; 4 loads. Sequential vs.\ pipelined
time and speedup?
\end{exercise}
\begin{solution}
Sequential: $4\times90 = 360$~min $=6$~h. Pipelined: the 40-min dryer sets
the pace: $30 + 4\times40 + 20 = 210$~min $=3.5$~h. Speedup $=6/3.5 =
1.7$ (not 3), because stages are unbalanced and the pipeline must fill and
drain.
\end{solution}

\begin{exercise}[Single-cycle vs.\ pipelined]
Stage times IF 2, Reg 1, ALU 2, Mem 2, WB 1~ns. Compare 3 loads, then 1003
loads, on a single-cycle (8~ns) and a pipelined (2~ns) machine.
\end{exercise}
\begin{solution}
3 loads: single-cycle $3\times8 = 24$~ns; pipelined $(5+2)\times2 = 14$~ns;
speedup $1.7$.
1003 loads: single-cycle $1003\times8 = 8024$~ns; pipelined
$(5+1002)\times2 = 2014$~ns; speedup $3.98 \approx 8/2 = 4$. The speedup
approaches the ratio of clock periods (not 5, because the stages are
unbalanced: the 1-ns stages still take a 2-ns cycle).
\end{solution}

\begin{exercise}[Pipelining a multi-cycle CPU]
Unpipelined: 1~ns clock, ALU/branch 4 cycles, memory 5 cycles; mix 40\%
ALU, 20\% branch, 40\% memory. Pipelining adds 0.2~ns to the clock.
Speedup?
\end{exercise}
\begin{solution}
Unpipelined average $=1\times(0.6\times4 + 0.4\times5) = 4.4$~ns.
Pipelined average $= 1.2$~ns (ideal $\CPI=1$). Speedup $=4.4/1.2 = 3.7$.
\end{solution}

\begin{exercise}[Structural hazard]
Loads/stores are 40\% of instructions; the machine with one memory port
stalls 1 cycle per data access but has a $1.05\times$ faster clock. Which is
better?
\end{exercise}
\begin{solution}
Hazard machine: $(1+0.4\times1)\times t/1.05 = 1.33\,t$ versus $1.0\,t$ for
the ideal machine, so the machine without the hazard is $1.33\times$
faster.
\end{solution}

\begin{exercise}[RAW hazard on R1]
\texttt{add r1,r2,r3; sub r4,r1,r3; and r6,r1,r7; or r8,r1,r9;
xor r10,r1,r11}. Stalls and total cycles without and with forwarding?
\end{exercise}
\begin{solution}
Without forwarding \texttt{sub} ($d=1$) needs 2 stalls; after that every
later reader is far enough behind:
\begin{lstlisting}[basicstyle=\ttfamily\footnotesize,keywords={}]
                c1    c2    c3    c4    c5    c6    c7    c8    c9    c10   c11
add r1,r2,r3    IF    ID    EX    MEM   WB
sub r4,r1,r3          IF    ID    stall stall EX    MEM   WB
and r6,r1,r7                IF    stall stall ID    EX    MEM   WB
or  r8,r1,r9                                  IF    ID    EX    MEM   WB
xor r10,r1,r11                                      IF    ID    EX    MEM   WB
\end{lstlisting}
11 cycles. With forwarding (EX/MEM $\to$ EX to \texttt{sub}, MEM/WB $\to$ EX
to \texttt{and}, register file for \texttt{or}/\texttt{xor}): 0 stalls,
$5+4 = 9$ cycles.
\end{solution}

\begin{exercise}[Load-use hazard]
\texttt{lw r1,0(r2); sub r4,r1,r6; and r6,r1,r7; or r8,r1,r9} with
forwarding.
\end{exercise}
\begin{solution}
The loaded value exists only at the end of MEM (cycle 4) but \texttt{sub}
needs it at the start of its EX (cycle 4): 1 stall is unavoidable.
\begin{lstlisting}[basicstyle=\ttfamily\footnotesize,keywords={}]
              c1    c2    c3    c4    c5    c6    c7    c8    c9
lw  r1,0(r2)  IF    ID    EX    MEM   WB
sub r4,r1,r6        IF    ID    stall EX    MEM   WB
and r6,r1,r7              IF    stall ID    EX    MEM   WB
or  r8,r1,r9                          IF    ID    EX    MEM   WB
\end{lstlisting}
9 cycles ($5+3+1$).
\end{solution}

\begin{exercise}[Compiler scheduling]
\texttt{a=b+c; d=e-f;} compiled as \texttt{lw Rb; lw Rc; add Ra; sw Ra;
lw Re; lw Rf; sub Rd; sw Rd} (forwarding, 1-cycle load delay). Stalls
before and after scheduling?
\end{exercise}
\begin{solution}
Original: \texttt{add} right after \texttt{lw Rc} and \texttt{sub} right
after \texttt{lw Rf} $\Rightarrow$ 2 stalls, $5+7+2 = 14$ cycles.
Scheduled \texttt{lw Rb; lw Rc; lw Re; add Ra; lw Rf; sw Ra; sub Rd; sw Rd}:
no load is immediately followed by its use $\Rightarrow$ 0 stalls, 12
cycles.
\end{solution}

\begin{exercise}[Branch stall impact]
30\% of instructions are branches with a 3-cycle stall. New CPI?
\end{exercise}
\begin{solution}
$\CPI = 1 + 0.30\times3 = 1.9$ -- almost twice the ideal. With the branch
resolved in ID (1 cycle): $1+0.3 = 1.3$.
\end{solution}

\begin{exercise}[Comparing branch schemes]
4\% unconditional, 6\% conditional untaken, 10\% conditional taken;
pipeline depth 5. Find CPI and speedups for stall, predict not taken,
predict taken, delayed branch (penalty 0.5).
\end{exercise}
\begin{solution}
\begin{tabular}[t]{@{}lcccc@{}}
Scheme & CPI & vs.\ unpipelined & vs.\ stall \\
\hline
Stall: $1+0.20\times3$ & 1.60 & $5/1.6=3.1$ & 1.00 \\
Not taken: $1+0.04\times1+0.10\times3$ & 1.34 & 3.7 & 1.19 \\
Taken: $1+0.14\times1+0.06\times2$ & 1.26 & 4.0 & 1.27 \\
Delayed: $1+0.20\times0.5$ & 1.10 & 4.5 & 1.45 \\
\end{tabular}


\end{solution}
\begin{remark}
\slidefix{} The slide lists 1.29 as the speedup of predict taken over stalling.
That is the ratio of the rounded speedups $4.0/3.1$; the exact value is
$1.6/1.26=1.27$.
\end{remark}

\begin{exercise}[CPI of a realistic mix]
Loads 30\% (25\% followed immediately by a use), branches 20\% (45\%
taken); forwarding, branch resolved in ID, predict not taken.
\end{exercise}
\begin{solution}
$\CPI = 1 + 0.3\times0.25\times1 + 0.2\times0.45\times1 = 1.165$.
\end{solution}

\begin{exercise}[Hazards, NOPs and forwarding]
\leavevmode
\begin{lstlisting}
(1) lw  $1, 40($6)
(2) add $6, $2, $2
(3) sw  $6, 50($1)
(4) add $4, $5, $6
(5) lw  $6, 10($4)
\end{lstlisting}
(A) Find all data dependencies. (B) Without forwarding, which cause hazards?
(C) Insert NOPs so the code runs correctly without forwarding. (D) How many
cycles does (C) take? (E) With full forwarding, how many cycles, and which
forwarding paths are used?
\end{exercise}
\begin{solution}
\textbf{(A)} RAW: (1)$\to$(3) on \texttt{\$1}; (2)$\to$(3) on \texttt{\$6};
(2)$\to$(4) on \texttt{\$6}; (4)$\to$(5) on \texttt{\$4}.
WAR: (1)$\to$(2), (3)$\to$(5), (4)$\to$(5) on \texttt{\$6}.
WAW: (2)$\to$(5) on \texttt{\$6}.

\textbf{(B)} Only RAW with $d\le2$: (1)$\to$(3) ($d=2$), (2)$\to$(3)
($d=1$), (2)$\to$(4) ($d=2$), (4)$\to$(5) ($d=1$). WAR/WAW are not hazards
here.

\textbf{(C)} Each consumer must be at least 3 slots after its producer:
\begin{lstlisting}
lw  $1, 40($6)
add $6, $2, $2
nop
nop
sw  $6, 50($1)
add $4, $5, $6
nop
nop
lw  $6, 10($4)
\end{lstlisting}
\textbf{(D)} 9 instructions $\Rightarrow 5 + 8 = 13$ cycles:
\begin{lstlisting}[basicstyle=\ttfamily\scriptsize,keywords={}]
               c1   c2   c3   c4   c5   c6   c7   c8   c9   c10  c11  c12  c13
lw  $1,40($6)  IF   ID   EX   MEM  WB
add $6,$2,$2        IF   ID   EX   MEM  WB
nop                      IF   ID   EX   MEM  WB
nop                           IF   ID   EX   MEM  WB
sw  $6,50($1)                      IF   ID   EX   MEM  WB
add $4,$5,$6                            IF   ID   EX   MEM  WB
nop                                          IF   ID   EX   MEM  WB
nop                                               IF   ID   EX   MEM  WB
lw  $6,10($4)                                          IF   ID   EX   MEM  WB
\end{lstlisting}

\textbf{(E)} No stalls are needed (the only load, (1), is used two
instructions later), so $5+4 = 9$ cycles:
\begin{lstlisting}[basicstyle=\ttfamily\footnotesize,keywords={}]
               c1    c2    c3    c4    c5    c6    c7    c8    c9
lw  $1,40($6)  IF    ID    EX    MEM   WB
add $6,$2,$2         IF    ID    EX    MEM   WB
sw  $6,50($1)              IF    ID    EX    MEM   WB
add $4,$5,$6                     IF    ID    EX    MEM   WB
lw  $6,10($4)                          IF    ID    EX    MEM   WB
\end{lstlisting}
Forwarding paths: cycle 5, MEM/WB $\to$ EX gives \texttt{\$1} (from (1)) and
EX/MEM $\to$ EX gives \texttt{\$6} (from (2)) to the \texttt{sw}; cycle 6,
MEM/WB $\to$ EX gives \texttt{\$6} (from (2)) to (4); cycle 7, EX/MEM $\to$
EX gives \texttt{\$4} (from (4)) to (5).
\end{solution}

% =============================================================================
\subsection{Cache}

\begin{exercise}[Address breakdown]
A 32~KiB cache has 64-byte blocks and 32-bit byte addresses. (a) Give the tag, index
and offset widths for direct-mapped, 4-way and fully associative organizations.
(b) Find the total storage of the 4-way write-back version. (c) Which set and tag
does \texttt{0x12345678} map to in the 4-way cache?
\end{exercise}
\begin{solution}
(a) Offset $\log_264=6$; 512 blocks. Direct mapped: index 9, tag 17. 4-way: 128 sets,
index 7, tag 19. Fully associative: tag 26. (b) Per block $512+19+1+1=533$ bits;
$512\times533=272{,}896$ bits $=34{,}112$ bytes. (c) Block address
\texttt{0x12345678}$/64=$ \texttt{0x48D159}; set $=$ \texttt{0x48D159} mod 128 $=89$;
tag $=$ \texttt{0x48D159}$/128=$ \texttt{0x91A2}.
\end{solution}

\begin{exercise}[Write policies]
A direct-mapped cache of 4 one-word blocks starts empty. The CPU performs
\texttt{W0 R0 W4 W4 R0 W8 R8 W0} (word addresses). Count the misses, blocks read from
memory and words written to memory for (a) write-through with no-write allocate,
(b) write-back with write allocate.
\end{exercise}
\begin{solution}
(a) Writes never fill the cache: W0 M, R0 M (fetch), W4 M, W4 M, R0 H, W8 M, R8 M
(fetch, replaces 0), W0 M: 7 misses, 2 blocks read, 5 words written. (b) W0 M (fetch,
dirty), R0 H, W4 M (write back 0, fetch 4), W4 H, R0 M (write back 4, fetch 0), W8 M
(0 clean, fetch 8), R8 H, W0 M (write back 8, fetch 0): 5 misses, 5 blocks read,
3 words written; block 0 is dirty at the end.
\end{solution}

\begin{exercise}[Multilevel cache and CPI]
L1 hit 1 cycle, miss rate 2\%; L2 hit 10 cycles, local miss rate 40\%; memory 100
cycles; 1.3 memory accesses per instruction; $\CPI_{exec}=1$. Find the AMAT and the
CPI with and without L2, and the L2 global miss rate.
\end{exercise}
\begin{solution}
With L2: $\AMAT=1+0.02\times(10+0.4\times100)=2$ cycles; stalls per instruction
$=1.3\times0.02\times50=1.3$, $\CPI=2.3$. Without: $\AMAT=1+0.02\times100=3$;
stalls $1.3\times0.02\times100=2.6$, $\CPI=3.6$. Speedup $3.6/2.3=1.57$. Global L2
miss rate $=0.02\times0.4=0.8\%$.
\end{solution}

% =============================================================================
\subsection{Virtual Memory}

\begin{exercise}[Address translation]
32-bit VA, 4~KiB pages; page table: VPN \texttt{00400} $\to$ PFN \texttt{0A1};
\texttt{00401} not present; \texttt{00403} $\to$ \texttt{1F2}. Translate
\texttt{0x00403A7C}, \texttt{0x00400008}, \texttt{0x00401010}.
\end{exercise}
\begin{solution}
The low 3 hex digits are the offset. \texttt{0x00403A7C} $\to$ \texttt{0x1F2A7C};
\texttt{0x00400008} $\to$ \texttt{0x0A1008}; \texttt{0x00401010}: page not present,
page fault.
\end{solution}

\begin{exercise}[Page table size]
(a) 32-bit VA, 4~KiB pages, 4-byte PTEs: flat table size per process? (b) Two levels
($10+10+12$): a process uses 8~MiB at the bottom and 1~MiB of stack at the top; how
much memory do its tables take? (c) 48-bit VA, 4~KiB pages, 8-byte PTEs: how many
levels if each fits in a page?
\end{exercise}
\begin{solution}
(a) $2^{20}\times4$~B $=4$~MiB. (b) Each second-level table maps 4~MiB: 2 for the code
and data, 1 for the stack, plus the root: $4\times4$~KiB $=16$~KiB. (c) VPN $=36$ bits;
a page holds $2^9$ PTEs, so $\lceil36/9\rceil=4$ levels.
\end{solution}

\begin{exercise}[TLB and page faults]
TLB 10~ns, memory 100~ns, TLB hit rate 95\%, two-level page table. (a) Effective
access time? (b) Add page faults on $10^{-6}$ of accesses costing 8~ms. (c) Which are
possible: TLB hit with page fault; TLB miss, page present, data in cache; TLB miss,
page fault, data in cache?
\end{exercise}
\begin{solution}
(a) $0.95\times110+0.05\times(10+3\times100)=104.5+15.5=120$~ns. (b)
$+10^{-6}\times8\times10^6=8$~ns: 128~ns. (c) Impossible; possible; impossible.
\end{solution}

\begin{exercise}[OPT and Clock]
3 frames, references \texttt{1 2 1 1 3 2 1 4 3 1 1 2 4 1 5 6 2 1}. Count the faults
with OPT and with Clock (new pages get bit 1; the pointer moves past a newly loaded
page).
\end{exercise}
\begin{solution}
OPT: faults at 1, 2, 3, 4 (evict 2), 2 (evict 3), 5 (evict 4), 6 (evict 5): 7 faults.
Clock: 1, 2, 3, 4 (all bits set: clear all, evict 1), 1 (evict 2), 2 (evict 3), 5
(evict 4), 6 (evict 1), 1 (evict 2): 9 faults. (FIFO gives 9, LRU 11.)
\end{solution}

% =============================================================================
\subsection{Multiprocessors}

\begin{exercise}[MSI vs.\ MESI]
Three processors, block $X$ in no cache. Give the states (P1 P2 P3) and bus actions
under MSI and MESI for: P1 reads; P2 reads; P1 writes; P3 reads; P2 writes; P1 reads.
Then P1 reads block $Y$ (cached nowhere) and writes it: how many bus transactions?
\end{exercise}
\begin{solution}
MSI: S I I (BusRd); S S I (BusRd); M I I (invalidate); S I S (BusRd, P1 writes back);
I M I (write miss); S S I (BusRd, P2 writes back). MESI is the same except the first
step gives E I I and the second changes P1 from E to S. For $Y$: MSI needs BusRd then an
invalidate (2 transactions); MESI loads $Y$ in E and upgrades to M silently (1).
\end{solution}

\begin{exercise}[Dependences and issue policies]
\texttt{I1: div r1,r2,r3}; \texttt{I2: add r4,r1,r5}; \texttt{I3: sub r5,r6,r7};
\texttt{I4: mul r4,r6,r8}; \texttt{I5: and r9,r5,r4}. List the dependences; say which
hazards matter with in-order issue and out-of-order completion, and with out-of-order
issue; rename to keep only true dependences.
\end{exercise}
\begin{solution}
RAW: I1$\to$I2 (r1), I3$\to$I5 (r5), I4$\to$I5 (r4). WAR: I2$\to$I3 (r5). WAW:
I2$\to$I4 (r4). Out-of-order completion: WAW (I4 finishing before I2). Out-of-order
issue: also WAR (I3 writing r5 before I2 reads it). Renamed: \texttt{I3: sub p5,r6,r7};
\texttt{I4: mul p4,r6,r8}; \texttt{I5: and r9,p5,p4}.
\end{solution}

\begin{exercise}[VLIW]
A 2-slot VLIW (one memory, one ALU slot); a load's result is usable 2 cycles after
issue, an ALU result 1 cycle after. Schedule \texttt{lw r1,0(r10)};
\texttt{lw r2,4(r10)}; \texttt{add r3,r1,r2}; \texttt{sw r3,8(r10)};
\texttt{addi r10,r10,12}.
\end{exercise}
\begin{solution}
Words: [\texttt{lw r1} | nop], [\texttt{lw r2} | nop], [nop | nop], [nop |
\texttt{add}], [\texttt{sw} | \texttt{addi}]: 5 words, 5 \texttt{nop}s (50\% of
slots). \texttt{sw} reads the old \texttt{r10} in the same word as \texttt{addi}.
\end{solution}

\begin{exercise}[Barrier reuse]
Why can the simple counter barrier deadlock when used in a loop, and how does sense
reversal fix it?
\end{exercise}
\begin{solution}
A fast thread leaves, reaches the barrier again and resets \texttt{release} before a
slow thread has seen it; the slow thread spins forever and the count never reaches
the total. With sense reversal each episode waits for the opposite value of
\texttt{release}, so a new episode never undoes the previous release.
\end{solution}

